\subsection{Further large-corner reductions at eighteen arrows}

\begin{lemma}[Proved: a trivial action at one end of a cut]
\label{CP2:A-trivial-cut-H2}
Let a category be a collage of a full source block $D$ and a
finite monoid $M$. Suppose $M$ acts trivially on every connecting
hom-set. Replacing $M$ by a trivial corner is an actual category
retract, and its pullback is an isomorphism on rational $H^2$.
The opposite statement also holds.
\end{lemma}
\begin{proof}
Let $T$ be the classifying space of the connecting $D$-act.
The two actual bar pushouts are
\[
 BD\longleftarrow T\times BM\longrightarrow BM,
 \qquad BD\longleftarrow T\longrightarrow *.
\]
The quotient functor has the section that includes the identity
of $M$. Every finite monoid has finite universal group, and
therefore $H^1(BM;\mathbb Q)=0$. Kunneth identifies the first
cohomology of $T\times BM$ with $H^1(T)$ and its second
cohomology with
\[
 H^2(T)\ \oplus\ H^0(T)\otimes H^2(BM).
\]
In the Mayer--Vietoris map from $H^2(BD)\oplus H^2(BM)$,
the two components are respectively the pullback to $T$ and
the injective map $b\mapsto-1_T\otimes b$.
Thus both the kernel in degree two and the preceding cokernel
in degree one are unchanged by the quotient. Naturality of
the exact sequence gives the stated $H^2$ isomorphism.
Here $T$ need only be nonempty; connectedness is not assumed.
\end{proof}

\begin{lemma}[Proved: a permutation orbit opposite a zero second group]
\label{CP2:A-permutation-fiber-cups}
In a collage $D\to M$, suppose $H^2(BD;\mathbb Q)=0$ and
one connecting stalk contains a nonempty finite $M$-invariant
subset on which every actor is a permutation. All rational
degree-two products of the collage vanish.
\end{lemma}
\begin{proof}
The homotopy-orbit space of that subset maps to the two-sided
bridge. Its projection to $BM$ is a finite covering, since
the action factors through a finite permutation group. The
covering transfer makes pullback from $H^2(BM;\mathbb Q)$
injective. Compatibility of the two restrictions of a collage
class therefore forces its restriction to $BM$ to be zero.
The restriction to $BD$ is zero by hypothesis. The two-block
relative suspension argument kills all second products.
\end{proof}

\begin{proposition}[Proved: two small incoming stalks through order eleven]
\label{CP2:A-two-small-incoming-eleven}
Suppose a full-linear category has first two corner orders
at most two, last corner order at most eleven, two arrows
from its middle to its last object, and at most three arrows
from its first to its last. If its rational homology is
bounded, every rational degree-two product vanishes. There
is no bound on its first adjacent hom-set.
\end{proposition}
\begin{proof}
First let the first two corners be trivial. Write $X,Y,Z$
for the three forward hom-sets, so $|Y|=2$, $|Z|\leq3$.
Every $x\in X$ supplies an $M$-equivariant map
$\beta_x:Y\to Z$. A constant map gives a fixed incoming
arrow. Otherwise every map is injective. Two different
two-point images in a three-set have a singleton invariant
intersection, again giving a fixed point. Hence all images
are one common invariant two-set $U$.

If the action on $Y$ consists of permutations, the preceding
lemma applies because the first two objects form a graph.
If two maps $Y\to U$ differ, their quotient is the
transposition of $Y$, commuting with every actor. A map of
two points commuting with their transposition is a
permutation, so this is the same case. Otherwise all
$\beta_x$ agree and some actor on $Y$ is constant. Its
action category is connected.

Fold $X$ to one chosen member, leaving $Y,Z,M$ fixed.
This is an actual category retract $C\to C_0$ because the
composition maps agree. The relative normalized chains,
one summand for each removed member of $X$, are the
augmentation cone of the $Y$ action-category chains,
shifted once. Their second homology is
$\widetilde H_0(Y_{hM};\mathbb Q)=0$. The retraction
therefore induces an isomorphism on $H^2$. The retract
has bounded homology. Deleting its middle object gives
the ordinary $M$-act cone on $Z$, and
Theorem~\ref{CP2:A-bounded-eleven-acts} kills its products.

For an initial corner of order at most two, take its
actual exact-orbit quotient first. This preserves the
rational ring, does not increase $|X|,|Z|$, and leaves
$Y$ unchanged. The initial corner is now trivial.
Let $N$ be the remaining middle corner, of order at most
two. Its rational algebra is semisimple.

If the right $N$-action on $Y$ is nontrivial, its orbit
set is a singleton. The equivariant graph model of
Lemma~\ref{CP2:A-baseline-semisimple-pair-graph} has
$Y/N$ as one vertex set. Its single vertex is fixed by
$M$, giving a section to the graph's homotopy-orbit
projection. The source block has zero rational $H^2$,
and the same two-block argument applies.

If the right action on $Y$ is trivial, replace $N$ by
a point and $X$ by $X/N$. Composition descends since
$\beta(y,nx)=\beta(yn,x)=\beta(y,x)$. This is an
actual rational homology equivalence. Indeed, the
strict-object bar filtration has only its corner,
edge and triangle terms. At $N$, exact semisimple
coinvariants identify respectively the corner bar
with a point, the $X$ bar with $X/N$, and the middle
triangle tensor with
$\mathbb QY\otimes\mathbb Q(X/N)$, since $N$ acts
trivially on $Y$. The other terms are unchanged.
The finite filtration proves the comparison and
passes boundedness to the quotient. The already
proved trivial-corner case now applies.
\end{proof}

\begin{proposition}[Proved: two first arrows and four long arrows]
\label{CP2:A-two-two-four-incoming}
A full-linear category with trivial first two corners,
last corner of order at most eleven, $|X|=|Y|=2$ and
$|Z|\leq4$ cannot have type $\CP^2$.
\end{proposition}
\begin{proof}
Consider its two maps $\beta_0,\beta_1:Y\to Z$.
Constant maps, distinct images meeting in one point,
or a permutation action on $Y$ are excluded by the
fixed-arrow or permutation-fiber lemmas. If the images
agree and the maps differ, the commuting transposition
again forces the permutation case. If the maps agree,
folding $X$ to a singleton is a cut-profunctor retraction.
Lemma~\ref{CP2:A-cut-retract} would retain projective-plane
type on its image. Deleting the middle object there
gives a bounded ordinary act cone over a monoid of
order at most eleven, a contradiction.

The only remaining case has disjoint images. They
exhaust $Z$, since both have size two. Fold $X$ onto
$x_0$ and fold the second image onto the first by
$\beta_0\beta_1^{-1}$. The latter map is $M$-equivariant.
These two maps form an actual idempotent natural
endomorphism of the cut from object0 to objects1,2:
its composition equation is exactly
$\gamma_Z\beta_i=\beta_0$.
Its image has $|X|=1$, $|Z|=2$. The same cut-retract
lemma and middle deletion give an ordinary binary act
cone, whose second products vanish universally.
\end{proof}

\begin{lemma}[Proved: a trivial middle corner and a singleton long hom-set]
\label{CP2:A-trivial-middle-singleton-long}
If a full-linear category has a trivial middle corner
and a singleton long hom-set, every rational
degree-two product vanishes, with no restriction on
its two other corners or adjacent hom-sets.
\end{lemma}
\begin{proof}
Let $D$ be the first two objects and $M$ the last.
Choose an arrow $y$ from the middle to the last object.
The unique long arrow and $y$ give a natural
transformation from the inclusion of $D$ to the
constant functor at the last object: every two-cross
composite is the unique long arrow, the first corner
fixes it, and the middle corner is trivial.
Thus $BD\to BC$ is null-homotopic. The unique long
arrow also gives a natural transformation from the
constant functor at the first object to the inclusion
of $M$, so $BM\to BC$ is null-homotopic.
Both restrictions of every positive class vanish.
Apply the two-block relative suspension argument.
\end{proof}

\begin{lemma}[Proved: a three-element middle opposite a binary stalk]
\label{CP2:A-small-middle-binary-singleton-long}
Suppose the first corner is trivial, the middle corner
$N$ has at most three elements, $|Y|=2$ and $|Z|=1$.
Every rational degree-two product vanishes.
\end{lemma}
\begin{proof}
The singleton $Z$ makes the last corner's inclusion
null-homotopic. If the right $N$-action on $Y$ has a
common fixed point, that point and $Z$ give a natural
transformation from the source block to the constant
last object, as above, and both restrictions vanish.

Suppose that action has no common fixed point. Its
transformation image has at most three elements. On
a two-set the possibilities are the permutation
group $C_2$ and the identity together with both
constants. In the former case $N$ surjects onto
$C_2$. Its minimal group has at least two elements;
the order bound forces $N=C_2$ or $C_2$ with a new
identity. Both rational algebras are semisimple, so
the source two-object block has zero rational $H^2$.
Together with the already null last-corner inclusion
this kills the products.

In the constants case the commuting last-corner
action on $Y$ is the identity: a map commuting with
both constants must fix both points. That corner
also acts trivially on $Z$. Apply
Lemma~\ref{CP2:A-trivial-cut-H2} to replace it by a
point. The resulting two trivial endpoint corners
and binary adjacent stalk satisfy the binary
projective-cover lemma, so all its second products
vanish. The $H^2$ isomorphism pulls back the result.
\end{proof}

\begin{proposition}[Proved: a three-element initial singleton]
\label{CP2:A-three-initial-singleton}
A full-linear category with initial corner order at
most three, trivial middle corner, last corner order
at most eleven, $|X|=1$, $|Y|=2$ and $|Z|\leq3$
cannot have type $\CP^2$.
\end{proposition}
\begin{proof}
If the initial trivial left module is projective,
use Proposition~\ref{CP2:A-eleven-projective-singleton}.
The small-monoid classification leaves only the
identity-adjoined two-element left-zero band.
Its two right actors $P,Q$ on $Z$ are retractions
onto a common image $U$, and the composition map
$\beta:Y\to Z$ lands in $U$. A constant $\beta$,
or $|U|=1$, supplies a fixed incoming arrow. The
source block has a terminal middle object, so its
second homology is zero and this case is excluded.
Otherwise $\beta$ is injective and $|U|\geq2$.

If $P=Q$, their common map on $Z$, together with
the identity on $X$, is an idempotent natural
endomorphism of the cut from the first object.
Its image retains a putative projective-plane type
by the cut-retract lemma. The initial corner acts
trivially on both image stalks, so the opposite
version of Lemma~\ref{CP2:A-trivial-cut-H2} replaces
it by a point without changing $H^2$. The replacement
is itself a category retract and hence bounded.
Middle deletion and the bounded-eleven act theorem
give the contradiction.

If $P\ne Q$, the size bound forces $|Z|=3$,
$|U|=2$, and one state outside $U$. The equivariant
graph of the right left-zero action has vertices
$U$ and the one edge from $Pz$ to $Qz$. Its two
endpoints are distinct. The element category of
the full incoming source-block act adds the two
vertices $Y$ and the edges from $\beta(y)$ to $y$.
This follows from the natural action-category graph
comparison: its fixed core vertices are unchanged,
and the unique first-adjacent arrow attaches exactly
the indicated two edges. All comparisons commute
with the last-corner action. The resulting graph
is a tree, hence contractible. The source block is
also contractible, so the two-block bar pushout is
the cofiber of an equivariant contractible graph's
homotopy-orbit map to the last corner. That map is
a homotopy equivalence, giving a contradiction.
\end{proof}

\begin{proposition}[Proved: three middle-to-last arrows and two long arrows]
\label{CP2:A-three-two-incoming}
A full-linear category with trivial first two corners,
$|Y|=3$, and $|Z|=2$ cannot have type $\CP^2$,
for any last monoid and any nonempty $X$.
\end{proposition}
\begin{proof}
Any constant composition map $\beta_x:Y\to Z$ gives
a fixed point. Assume they are all surjective. If two
have the same partition and differ, they differ by
the transposition of $Z$, so every last-corner actor
on $Z$ is a permutation. The permutation-fiber lemma
excludes this case. If all maps agree, fold $X$ to
a singleton by a cut retraction. Its image would
still be a projective plane, but middle deletion
gives an ordinary two-state cone, a contradiction.

Thus choose two maps with different two-part
partitions. Their joint map embeds the three-set
$Y$ into $Z^2$. The diagonal transposition of $Z^2$
has no fixed point and cannot preserve a three-set.
No last-corner actor can therefore transpose $Z$.
Unless there is a globally fixed point, the actor
image on $Z$ contains both constants. The joint
image consequently consists of $(0,0),(1,1)$ and
one off-diagonal pair. Name the corresponding states
$y_0,y_1,w$, and interchange the two chosen maps
if needed so that the last pair is $(0,1)$.
An actor is the identity on $Y$ when it is the
identity on $Z$, and is constant to $y_i$ when
it is constant to $i$. Every other $\beta_x$ must
take $y_i$ to $i$ and $w$ to either0 or1. Thus
there are exactly two composition maps in use.

Fold all repetitions in $X$, retaining one copy of
each map. This is a cut retraction fixing $Y,Z$.
Its image would still have projective-plane type.
There $|X|=2$ and the first two objects form a
circle $BD$. The incoming act graph has vertex
sets $Y,Z$ and the six edges of the two maps.
Contract the three edges of the first map. They
form an equivariant forest. The quotient graph
$\Gamma$ has two vertices, one loop at each,
and a connecting edge. It is connected and its
first homology is the permutation module
$\mathbb Q\{0,1\}$: the two constant actors send
both loop generators to the selected loop.

Put $X_0=\operatorname{hocofib}(\Gamma_{hM}\to BM)$.
For a connected graph its augmented chain complex
is equivariantly quasi-isomorphic to $H_1(\Gamma)[1]$.
Hence
\[
 H_2(X_0;\mathbb Q)=H_1(\Gamma;\mathbb Q)_M=\mathbb Q.
\]
Collapsing $BD$ in the collage gives $X_0$.
Its cofiber boundary $H_2(X_0)\to H_1(BD)$ sends
each loop generator to the generator of $BD$:
before contraction it follows the second first-arrow
and returns along the first. It is therefore the
augmentation on this permutation module, an
isomorphism after coinvariants. Since $H_2(BD)=0$,
the cofiber exact sequence forces $H_2(BC)=0$.
The retract cannot be a projective plane.
\end{proof}

\begin{proposition}[Proved: a two-element endpoint in the three-by-three triangle]
\label{CP2:A-three-three-two-endpoint}
Every full-linear category with profile
$(1,m,2;3,3,1)$ and $m+9\leq22$ is excluded.
\end{proposition}
\begin{proof}
Take the exact-orbit quotient at the last corner,
whose rational algebra is semisimple. It preserves
the rational cohomology ring. If its three-element
incoming stalk shrinks, it becomes binary and the
binary projective-cover obstruction applies. If
it does not shrink, the last corner acts trivially
on that stalk and on the singleton long stalk.
The quotient then has an actual section. Its integral
homology groups are direct summands of those of a
putative projective plane. The rational $H^2$
isomorphism forces its integral $H_2$ to be $\mathbb Z$;
naturality of the primitive square then forces
$H_4=\mathbb Z$ with the same unit square. Its
fundamental group is a retract of the trivial group.
The projective-plane recognition argument of the
cut-retract lemma therefore shows that this quotient
also has type $\CP^2$.
It has trivial endpoints and singleton long hom-set.
Triangle compression gives a monoid of order $m+9$,
contradicting the monoid minimum of twenty-three.
\end{proof}

\begin{remark*}[Exact remaining large-corner scope]
The fresh eighteen-arrow profile ledger initially has fourteen
full-linear profiles with a corner of order at least eight.
The preceding lemmas remove twelve; the full three-by-three
triangle theorem removes $(1,8,1;3,3,2)$.
The single remaining profile is $(1,1,8;2,3,3)$, together
with its opposite. Its separate complete reduction and fresh
literal calculation are given in
Theorem~\ref{CP2:A-extreme-eight-two-three-three} below.
\end{remark*}
